Полностью нормализованная реляционная модель строится на основе
PostgreSQL JSONB-модели. Основные сущности, связи, суррогатные первичные
ключи, альтернативные ключи и внешние ключи сохраняются. Заметно изменяется
способ хранения составных и вариативных данных:
поля, которые в JSONB-модели находились внутри \texttt{jsonb}, раскрываются
в столбцы и отдельные отношения.

\subsubsection{Нормализация справочных значений}

Сначала из отношений выделяются домены допустимых
значений. Тип территории хранится в \texttt{area\_type}, тип зоны --- в
\texttt{zone\_type}, состояние камеры --- в \texttt{camera\_status}, тип
события --- в \texttt{event\_type}, уровень значимости события --- в
\texttt{event\_severity}. По тому же правилу выделяются
\texttt{object\_type}, \texttt{direction}, \texttt{video\_codec} и
\texttt{signal\_lost\_reason}.

Такие таблицы выделяются при устранении транзитивных зависимостей, что
соответствует приведению схемы к третьей нормальной форме. Сам по себе код
статуса или типа не нарушает нормальную форму. Проблема возникает, если в
основной таблице вместе с кодом хранить его название, описание или ранг: эти
поля зависят уже не от ключа основной таблицы, а от справочного кода.
Например, идентификатор камеры определяет статус, а статус определяет его
название. Поэтому справочные атрибуты переносятся в отдельные отношения, а в
основных таблицах остаются внешние ключи. Такие отношения также задают
допустимые значения полей.

\subsubsection{Декомпозиция составных атрибутов}

Дальше раскладываются составные и многозначные атрибуты:

\begin{itemize}
  \item \texttt{camera.position} не образует отдельного отношения, так как
        описывает одну камеру и содержит фиксированный набор значений. Поэтому
        координаты установки и углы ориентации переносятся в столбцы таблицы
        \texttt{camera};

  \item \texttt{camera.settings} содержит несколько логически разных групп.
        Параметры видеопотока переходят в
        \texttt{camera\_stream\_setting}, настройки аналитики --- в
        \texttt{camera\_analytics\_setting}. Области детекции становятся
        отношением \texttt{camera\_detection\_zone}, а точки контура этих
        областей --- отношением \texttt{camera\_detection\_zone\_point}.
        Линии пересечения выносятся в \texttt{camera\_crossing\_line};

  \item \texttt{event.payload} зависит от типа события, поэтому общая часть
        события остается в \texttt{event}, а параметры, зависимые от типа события,
        переносятся в отношения \texttt{motion\_event\_detail},
        \texttt{object\_detection\_event\_detail},
        \texttt{line\_crossing\_event\_detail} и
        \texttt{signal\_lost\_event\_detail};

  \item массив \texttt{event.payload.objects} является многозначным
        атрибутом события. Его элементы хранятся в
        \texttt{detected\_object}. Атрибуты, зависящие от типа объекта,
        вынесены в \texttt{person\_object\_attribute} и
        \texttt{vehicle\_object\_attribute};

  \item \texttt{camera\_telemetry.metrics} содержит фиксированный набор
        технических показателей, поэтому раскрывается в столбцы таблицы
        \texttt{camera\_telemetry}: температуру, загрузку процессора,
        использование памяти, битрейт, потери пакетов, задержку и время
        непрерывной работы.
\end{itemize}

\subsubsection{Нормальная форма}

Декомпозиция выполняется ради устранения зависимостей, которые затрудняют
обновление данных и приводят к их дублированию. На первом шаге схема
приводится к первой нормальной форме: атрибуты становятся атомарными
относительно выбранной реляционной модели, поэтому
\texttt{event.payload.objects} и точки областей детекции выносятся в
отдельные отношения.

Затем проверяется зависимость неключевых атрибутов от ключей. Для отношений с
составными ключами неключевые атрибуты должны зависеть от всего ключа, а не от
его части. Например, в таблице связи \texttt{event\_camera} нет собственных
описательных атрибутов камеры или события: они остаются в \texttt{camera} и
\texttt{event}. Это соответствует требованию второй нормальной формы.

Третья нормальная форма требует, чтобы неключевые атрибуты не зависели от
других неключевых атрибутов. Поэтому названия, описания и ранги справочных
значений вынесены из основных таблиц в \texttt{area\_type},
\texttt{camera\_status}, \texttt{event\_type}, \texttt{event\_severity} и
другие справочники. В основной таблице остается только ссылка на справочное
значение.

С учетом заданных ключей и выделенных таблиц отношения могут
рассматриваться как удовлетворяющие более строгому условию BCNF: в каждой
нетривиальной функциональной зависимости детерминант является суперключом
своего отношения. Иначе говоря, атрибуты в каждой таблице описывают только
объект, заданный ее ключом, а данные о других объектах вынесены в собственные
отношения.


\begin{figure}
    \centering
    \includegraphics[
        width=\textwidth,
        height=1\textheight,
        keepaspectratio
    ]{rel_norm.pdf}
    \caption{Нормализованная реляционная модель}
\end{figure}

\FloatBarrier
